\documentclass{ximera}

\title{This is the sample!}

\begin{document}

\begin{abstract}
HELLO again and $\arcsin$.
\end{abstract}

\maketitle

\begin{tikzpicture}
  \draw (0,0) circle (1in);
  \draw (0,0) circle (0.25in);
  \draw (0,0) -- (1in,1in);
\end{tikzpicture}

 \begin{problem}
   The tolerance (0.01) means $8 \approx \answer[tolerance=0.01]{8}$
 \end{problem}

 \begin{problem}
   The tolerance (0.01) means $\pi \approx \answer[tolerance=0.01]{3.141592653}$
 \end{problem}

\includeinteractive{example.js}


\begin{foldable}
  So this should be foldable.

  And even this $2^x$ thing should be.
  
  \begin{foldable}
  But there is a subset.  \unfoldable{BUT NOT HERE.}
  \end{foldable}

  \unfoldable{But not this.}  That is a key point.  \unfoldable{So is this.}

 \begin{problem}
   The tolerance (17) means $3421 \approx \answer[tolerance=17]{3421}$
 \end{problem}

\end{foldable}

%% \begin{code}
%%   var Expression = require('math-expressions');

%%   Expression.fromText('x^2');
%% \end{code}

%% \begin{verbatim}
%%   Hello!  This is verbatim.
%% \end{verbatim}

%% \begin{problem}
%%   We can add such as in $2 + 2 = \answer{4}$.
  
%%   \begin{problem}
%%     Multiplication looks like $3 \times 3 = \answer{9}$.
    
%%     \begin{problem}
%%       Now consider $\sqrt{\answer{16}} = 4$.
      
%%       \begin{problem}
%%         Therefore $4 \times 4 = \answer{16}$.
%%       \end{problem}
%%     \end{problem}
%%   \end{problem}
%% \end{problem}

%% \begin{theorem}
%% Look at my shadow!
%% \end{theorem}

%% \begin{problem}
%%   \begin{multipleChoice}
%%     \choice{Incorrect}
%%     \choice{Wrong}
%%     \choice[correct]{It's this one}
%%     \choice{Not Right}
%%   \end{multipleChoice}

%%   \begin{problem}
%%     There were $\answer{4}$ possible answers to that question.

%%     \begin{problem}
%%       \begin{multipleChoice}
%%         \choice{Not correct}
%%         \choice[correct]{Pick me!}
%%         \choice{False}
%%         \choice{Untrue}
%%       \end{multipleChoice}
%%     \end{problem}
%%   \end{problem}
%% \end{problem}

\end{document}

%%% Local Variables:
%%% mode: latex
%%% TeX-master: t
%%% End:
% Changed at Mon Jul 27 22:23:04 EDT 2015
% Changed at Mon Jul 27 22:24:20 EDT 2015
% Changed at Mon Jul 27 22:25:23 EDT 2015
% Changed at Mon Jul 27 22:30:13 EDT 2015
% Changed at Mon Jul 27 22:32:16 EDT 2015
% Changed at Mon Jul 27 22:33:48 EDT 2015
% Changed at Tue Jul 28 00:05:31 EDT 2015
% Changed at Tue Jul 28 00:05:33 EDT 2015
% Changed at Tue Jul 28 00:05:34 EDT 2015
% Changed at Tue Jul 28 00:05:35 EDT 2015
% Changed at Tue Jul 28 00:05:55 EDT 2015
% Changed at Tue Jul 28 00:05:56 EDT 2015
% Changed at Tue Jul 28 00:05:57 EDT 2015
% Changed at Tue Jul 28 00:05:58 EDT 2015
% Changed at Tue Jul 28 00:05:59 EDT 2015
% Changed at Tue Jul 28 00:06:00 EDT 2015
% Changed at Tue Jul 28 00:06:01 EDT 2015
% Changed at Tue Jul 28 00:06:03 EDT 2015
% Changed at Tue Jul 28 00:06:04 EDT 2015
% Changed at Tue Jul 28 00:06:28 EDT 2015
